\documentclass[11pt]{article}
\usepackage[margin=1.25in]{geometry}
\usepackage{microtype}
\pagestyle{empty}

\begin{document}

\noindent
Peter Cotton\\
\texttt{peter.cotton@gsmc.ai}\\[2em]

\noindent
The Editors\\
\emph{Economics Letters}\\[2em]

\noindent
Dear Editors,

\medskip

Please consider the enclosed manuscript, ``Winning the large language capability battle and losing the production economy,'' for publication in \emph{Economics Letters}.

The paper has an unusual origin, which I mention because it explains both the question and the stakes. I operate roughly 70,000 private agentic prediction markets for a hedge fund. We decided to use them to assess the economic competence of large language models in a new way, and we guessed that profit divided by token cost might be closer to the mark than most published metrics. This paper is the result of trying to sharpen that guess.

The sharpening produced more than I expected. The note gives a closed-form competitive equilibrium in which a model strictly more accurate than every rival on every task nonetheless loses its entire share of production compute as tasks subdivide, because allocation runs on value per token rather than on capability. It also gives an identity: task-level average PnL per token equals the shadow price of compute divided by the elasticity of task value with respect to effective compute, so the guessed statistic ranks productivity exactly when that elasticity declines with crowding, and can mislead otherwise. Concavity alone is not sufficient, and a smooth counterexample within the stated assumptions is included.

Only afterwards did I appreciate what this means for the present contest between China and the United States. Export controls, national-scale capital, and policy attention are being steered by capability benchmarks. It is remarkable that the economically wrong measurement is in ubiquitous use; the paper states the right one, and the conditions under which it can actually be computed.

For transparency: I operate \texttt{modelrank.net}, the realized-PnL ranking discussed in Section~4 of the manuscript; this is disclosed in the manuscript's declarations. The manuscript is not under consideration elsewhere, and no data were used.

\medskip

\noindent
Sincerely,\\[1em]
Peter Cotton

\end{document}
